\subsection{Sprint 0 --- 11/08/2023 até 25/08/2023}

A Sprint 0 teve como objetivo preparar a equipe para as próximas etapas do projeto. Nessa fase, o foco principal não estava no desenvolvimento intensivo da aplicação, mas na definição das tecnologias, das \textit{user stories}, dos fluxos de trabalho, dos primeiros mockups e das bases técnicas que sustentariam o Plan Line.

O início da sprint foi marcado pela apresentação da equipe, o que ajudou a construir uma base de colaboração para o restante do semestre. Em seguida, passamos a discutir as estruturas fundamentais do projeto, incluindo a escolha das tecnologias, a configuração inicial do banco de dados, o levantamento de requisitos e a organização das primeiras histórias de usuário.

Como grande parte das tecnologias escolhidas ainda estava além do meu conhecimento naquele momento, dediquei uma parcela significativa da sprint ao estudo. Busquei me familiarizar com HTML, CSS, JavaScript e os primeiros conceitos de TypeScript. Esse período foi desafiador, pois evidenciou minha inexperiência em desenvolvimento de software, mas também foi muito enriquecedor por me colocar em contato direto com ferramentas utilizadas em um projeto real.

Durante o levantamento de requisitos, a equipe utilizou técnicas como \textit{brainstorming}, \textit{t-shirt sizing} e \textit{planning poker}. Essas práticas ajudaram na pontuação das histórias, na identificação de riscos e na discussão sobre as expectativas dos stakeholders. Para mim, foi importante perceber como a definição de requisitos depende não apenas de conhecimento técnico, mas também de comunicação, escuta e capacidade de transformar ideias amplas em tarefas executáveis.

O encerramento da sprint incluiu o primeiro encontro com os stakeholders Cassio, Letícia e Augusto, representantes da Triider. A reunião foi produtiva, mas também revelou dificuldades de clareza em alguns requisitos, especialmente relacionados à parte web da aplicação. Essa lacuna impactou o andamento inicial do projeto e mostrou a importância de validar entendimento com o cliente desde cedo.

Ao final da Sprint 0, realizamos uma retrospectiva e a equipe ficou satisfeita com o progresso inicial. Pessoalmente, essa sprint representou um choque de realidade. Eu estava entusiasmado por participar do meu primeiro projeto AGES, mas também percebi que precisaria lidar com uma curva de aprendizado intensa. Essa percepção foi importante para reforçar minha determinação em estudar, contribuir com a equipe e transformar a dificuldade inicial em oportunidade de crescimento.
